
Loading diagram...
Lightness:
0
Sync label/edge colours:
Preset:
LaTeX:(None)
Diagram:
Keyboard shortcutsctrl/

For a guide on using keyboard shortcuts in quiver see the tutorial.
General
Dismiss errors, and panels;
Cancel modification or movement;
Hide focus point;
Deselect, and dequeue cells	esc
Import from LaTeX	ctrlI
Export to LaTeX or Typst	ctrlE
Navigation
Pan view	Scroll
Enable mouse panning	ctrl,alt
Enable touch panning	Long press
Enable mouse zooming	⇧
Move focus point	←,↑,↓,→
Select next queued cell	⇥
Select previous queued cell	shift⇥
Select / deselect object	S
Select cells	;
Toggle cell selection	'
Modification
Focus / defocus label input	↵
Create object, and connect to selection	        
Move selected objects	B
Change source	,
Change target	.
Create arrows from selection	/
Copy	ctrlC
Cut	ctrlX
Paste	ctrlV
Styling
Reverse arrows	R
Flip arrows	E
Flip labels	F
Left-align labels	V
Centre-align labels	C
Over-align labels	X
Modify label position	I
Modify offset	O
Modify curve	K
Modify radius	N
Modify length	L (hold ⇧ to shorten symmetrically)
Modify level	M
Modify style	D
Display as arrow	A
Display as adjunction	J
Display as pullback/pushout	P (press again to switch corner style)
Modify label colour	Y
Modify arrow colour	U
Toolbar
Save diagram in URL	ctrlS
Undo	ctrlZ
Redo	shiftctrlZ
Select all	ctrlA
Select connected components	shiftctrlC
Deselect all	shiftctrlA
Delete	⌫,del
Centre view	G
Zoom out	ctrl-
Zoom in	ctrl=
Toggle grid	H
Toggle hints	shiftctrlH
Export
Toggle diagram centring	C
Toggle ampersand replacement	A
Toggle cramped spacing	R
Toggle standalone	S
Toggle fixed size	F
About

quiver is a modern, graphical editor for commutative and pasting diagrams, capable of rendering high-quality diagrams for screen viewing, and exporting to LaTeX via tikz-cd and Typst via fletcher.

Creating and modifying diagrams with quiver is orders of magnitude faster than writing the equivalent LaTeX or Typst by hand and, with a little experience, competes with pen-and-paper. To learn how to use quiver efficiently, see the tutorial.

The editor is open source and may be found on GitHub. If you would like to request a feature, or want to report an issue, you can do so here.

You can follow quiver on Mastodon or Twitter for updates on new features.
Thanks to

    S. C. Steenkamp, for helpful discussions regarding the aesthetic rendering of arrows.
    AndréC, for the custom TikZ style for curves of a fixed height.
    Andrew Stacey, for the custom TikZ style for shortened curves.
    Théophile Cailliau, for implementing Typst support.
    Nathan Corbyn, for adding the ability to export embeddable diagrams to HTML.
    Paolo Brasolin, for adding offline support.
    Carl Davidson, for discussing and prototyping loop rendering.
    Everyone who has improved quiver by submitting pull requests, reporting issues or suggesting improvements.

Created by varkor.
Welcome

quiver is a modern, graphical editor for commutative and pasting diagrams, capable of rendering high-quality diagrams for screen viewing, and exporting to LaTeX via tikz-cd and Typst via fletcher.

quiver is intended to be intuitive to use and easy to pick up. Here are a few tips to help you get started:

    Click and drag to create new arrows: the source and target objects will be created automatically.
    Double-click to create a new object.
    Edit labels with the input bar at the bottom of the screen.
    Click and drag the empty space around a object to move it around.
    Hold Shift (⇧) to select multiple cells to edit them simultaneously.

For a detailed guide to using quiver, see the tutorial.
Remember to include \usepackage{quiver} in your LaTeX preamble. You can install the package using CTAN, or open quiver.sty in a new tab to copy-and-paste.updated on 2025-07-05
% https://q.uiver.app/#q=WzAsMTAsWzIsMiwiWHhZIl0sWzQsMiwiWSJdLFsyLDQsIlgiXSxbNCw0LCJDIl0sWzEsMSwiUCJdLFszLDEsInQ9ZiBvIG0iXSxbMywzLCJnIG8gZj1ob2kiXSxbMSwzLCJ1PSBpb20iXSxbMiwwLCJQcm9waWVkYWQgVW5pdmVyc2FsIGRlbCBQdWxsYmFjayJdLFswLDAsIlEiXSxbNCwwLCIiLDAseyJzdHlsZSI6eyJib2R5Ijp7Im5hbWUiOiJkYXNoZWQifX19XSxbMCwxLCJmIl0sWzAsMiwiaSJdLFsyLDMsImgiLDJdLFsxLDMsImciLDJdLFs0LDIsIiIsMix7ImN1cnZlIjozfV0sWzQsMSwiIiwyLHsibGFiZWxfcG9zaXRpb24iOjYwLCJjdXJ2ZSI6LTJ9XSxbOSw0XSxbOSwyLCIiLDIseyJjdXJ2ZSI6NX1dLFs5LDEsIiIsMix7ImN1cnZlIjotNSwic3R5bGUiOnsiYm9keSI6eyJuYW1lIjoic3F1aWdnbHkifX19XV0=
\[\begin{tikzcd}
	Q && {Propiedad Universal del Pullback} \\
	& P && {t=f o m} \\
	&& XxY && Y \\
	& {u= iom} && {g o f=hoi} \\
	&& X && C
	\arrow[from=1-1, to=2-2]
	\arrow[curve={height=-30pt}, squiggly, from=1-1, to=3-5]
	\arrow[curve={height=30pt}, from=1-1, to=5-3]
	\arrow[dashed, from=2-2, to=3-3]
	\arrow[curve={height=-12pt}, from=2-2, to=3-5]
	\arrow[curve={height=18pt}, from=2-2, to=5-3]
	\arrow["f", from=3-3, to=3-5]
	\arrow["i", from=3-3, to=5-3]
	\arrow["g"', from=3-5, to=5-5]
	\arrow["h"', from=5-3, to=5-5]
\end{tikzcd}\]
